\chapter*{마무리와 다음 Phase 예고}
\addcontentsline{toc}{chapter}{마무리와 다음 Phase 예고}

Phase 0에서는 BERT 없이도 모델, 출력 헤드, 활성화 함수, Loss의 관계를 충분히 관찰했습니다. Phase 1에서는 같은 질문을 DistilBERT 위에서 다시 물었습니다. 모델은 커졌지만 질문은 그대로입니다. 출력은 몇 차원인가, sigmoid인가 softmax인가, Loss는 무엇을 벌점으로 보는가.

다음 Phase에서는 한국어 BERT로 같은 흐름을 압축해서 재방문합니다. 이 구조를 유지하면 새로운 모델을 만날 때도 덜 흔들립니다. 모델 이름보다 먼저 출력 모양과 라벨 형식을 확인하는 습관이 생기기 때문입니다.
